\section{Discussion}
\markright{Discussion}

In this experiment, various electronic components and laboratory instruments were studied, including resistors, capacitors, diodes, transistors, operational amplifiers, voltage regulators, multimeters, oscilloscopes, function generators, breadboards, veroboards, and soldering tools. Their construction, operating principles, specifications, and laboratory applications were examined.

The experiment highlighted the difference between passive components, which store or dissipate energy, and active devices, which control or amplify signals using an external power supply. It also demonstrated the oscilloscope's ability to analyze waveforms beyond the capabilities of a multimeter.

Proper component identification and handling were found to be essential. Correct resistor values, capacitor polarity, diode orientation, and the use of a current-limiting resistor with LEDs are critical for accurate measurements and preventing component damage.
